\section{Model framework}
\label{sec:model}

This section summarizes the one-dimensional framework of the companion
methods paper \cite{xue2026model} to the extent needed to interpret the
geysering computations, and then describes the vertical-shaft branch,
junction closures, and trapped-pocket treatment used in the geysering
configurations. The derivations, the Regime-Transition Riemann solver, and
the conservative cut-cell front-tracking algorithm are given in
\cite{xue2026model} and are not repeated here.

\subsection{Decoupled two-fluid mixed-flow formulation}
\label{sec:model:framework}

The framework describes a closed conduit as a union of liquid-full
(pressurized) and stratified gas--liquid reaches separated by moving regime
fronts. In stratified reaches the flow is governed by the decoupled
two-fluid balance law of \cite{xue2026model},
\begin{equation}
\frac{\partial \mathbf{U}}{\partial t}
+\frac{\partial \mathbf{F}}{\partial x}
=\mathbf{S},
\label{eq:system}
\end{equation}
with conservative state, flux, and regular source vectors
\begin{equation}
\mathbf{U}=
\begin{bmatrix}
A_g\rho_g \\ \rho_g Q_g \\ A_l \\ Q_l
\end{bmatrix},
\qquad
\mathbf{F}=
\begin{bmatrix}
\rho_g Q_g \\[2pt]
\rho_g Q_g^2/A_g+P_g A_g \\[2pt]
Q_l \\[2pt]
Q_l^2/A_l+\tfrac{1}{2}\Lambda_d A_l^2
\end{bmatrix},
\qquad
\mathbf{S}=
\begin{bmatrix}
0 \\[2pt]
P_g\dfrac{\partial A_g}{\partial x}
-A_g\rho_g g\cos\theta\dfrac{\partial h_l}{\partial x}+\mathcal{S}_g \\[2pt]
0 \\[2pt]
\mathcal{S}_d
\end{bmatrix},
\label{eq:vectors}
\end{equation}
where $A_g$ and $A_l$ are the gas and liquid cross-sectional areas, $\rho_g$
the gas density, $Q_g$ and $Q_l$ the gas and liquid volumetric flow rates,
$P_g$ the gas pressure, $h_l$ the liquid depth, $\theta$ the pipe
inclination, and $\mathcal{S}_g$, $\mathcal{S}_d$ the regular gas- and
liquid-momentum sources containing axial gravity, wall friction, and
interfacial shear; the shear closures follow standard stratified-flow
correlations \cite{issa2003,taitel1976}. Gas mass and gas momentum are
conserved variables, so gas transport, gas inertia, gas compressibility, and
shear-driven interfacial growth are part of the computed state. The
restoring coefficient in the liquid momentum flux is
\begin{equation}
\Lambda_d
=\underbrace{\frac{2gH_g}{A_l}}_{\text{gas pressure}}
+\underbrace{\frac{\rho_l-\rho_g}{\rho_l}\,g\zeta}_{\text{hydrostatic--buoyancy}}
-\underbrace{\frac{\rho_g(u_g-u_l)^2}{\rho_l A_g}}_{\text{IKH slip}},
\label{eq:lambda}
\end{equation}
which combines the gas-pressure, hydrostatic--buoyancy, and inviscid
Kelvin--Helmholtz (IKH) slip contributions; here $H_g$ is the gas pressure
head, $u_g$ and $u_l$ the phase velocities, and
$\zeta=\cos\theta/[D\sin(\gamma/2)]$ a circular-section geometric factor
with $D$ the pipe diameter and $\gamma$ the liquid central angle. The
eigenvalues of the decoupled liquid flux Jacobian,
\begin{equation}
\lambda^{\pm}=u_l\pm\sqrt{\Lambda_d A_l},
\label{eq:eigen}
\end{equation}
show that the liquid subsystem is strictly hyperbolic for $\Lambda_d>0$ and
reaches the neutral interfacial-instability (IKH) threshold at
$\Lambda_d=0$; beyond it, interfacial waves grow and liquid bridging can
occur.

In liquid-full reaches the gas phase is absent and the pressurized branch is
governed by the two-variable system
\begin{equation}
\frac{\partial A_l}{\partial t}+\frac{\partial Q_l}{\partial x}=0,
\qquad
\frac{\partial Q_l}{\partial t}
+\frac{\partial}{\partial x}\!\left(\frac{Q_l^2}{A_l}+gA_f h_s\right)
=\mathcal{S}_p,
\label{eq:pressurized}
\end{equation}
closed by the elastic-pipe relation
\begin{equation}
P-P_{\mathrm{ref}}=\rho_l a^2\,\frac{A_l-A_f}{A_f},
\qquad
h_s=\frac{P-P_{\mathrm{ref}}}{\rho_l g},
\label{eq:elastic}
\end{equation}
with $a$ the (numerical) pressure-wave speed, $A_f$ the undeformed pipe
area, $h_s$ the piezometric head, and $\mathcal{S}_p$ the regular
pressurized-branch source; this branch recovers the water-hammer limit. In
limiting cases the same area-variable form of
Eqs.~(\ref{eq:system})--(\ref{eq:elastic}) recovers water hammer,
open-channel gravity flow (with $\Lambda_d\to g\cos\theta/B_l$, where $B_l$
is the free-surface top width), interfacial-instability dynamics, and the
near-full mixed-flow balance \cite{xue2026model}.

Moving pressurized--stratified fronts are internal boundaries between the
two-variable system~(\ref{eq:pressurized}) and the four-variable
system~(\ref{eq:system}). Their local closure and the conservative
front-tracking update are summarized in
Sections~\ref{sec:model:rtrp}--\ref{sec:model:fv}, so that the computations
reported here can be interpreted and reproduced from the present paper
alone; the derivations, the construction of the active-set solver, and the
verification of the scheme on canonical benchmarks are given in
\cite{xue2026model}.

Two properties of this formulation matter for geysering. First, because gas
mass is conserved per computational region, entrapped pockets are identified
as connected gas regions bounded by liquid seals, and their pressure follows
from an isothermal equation of state $P_g = m_g R T / V_g$ evaluated with the
current pocket mass $m_g$ and volume $V_g$; no polytropic pocket law with a
prescribed volume history is imposed. Second, because the stratified branch
carries gas momentum, air intrusion into a shaft and counter-current
air--water exchange at the shaft mouth are resolved rather than prescribed.

\subsection{Regime-Transition Riemann closure at moving fronts}
\label{sec:model:rtrp}

At a moving pressurized--stratified front $x=x_\Gamma(t)$ with speed
$w_\Gamma=\dot{x}_\Gamma$, the local coordinate is oriented so that the
pressurized branch lies on the left and the stratified branch on the right;
the opposite case follows by reflection. Only liquid mass and liquid
momentum are common to both branches, and the ideal front has zero thickness
and stores neither. Integrating the mixed-regime liquid balance over a
control volume shrinking onto the front therefore gives the moving
Rankine--Hugoniot conditions
\begin{equation}
\begin{gathered}
Q_{p,\Gamma}-Q_{f,\Gamma}
=w_{\Gamma}\left(A_{p,\Gamma}-A_{l,\Gamma}\right),\\[2pt]
\left(\frac{Q_{p,\Gamma}^2}{A_{p,\Gamma}}+gA_f h_{s,\Gamma}\right)
-\left(\frac{Q_{f,\Gamma}^2}{A_{l,\Gamma}}
+\frac{1}{2}\Lambda_{d,\Gamma}A_{l,\Gamma}^2\right)
=w_{\Gamma}\left(Q_{p,\Gamma}-Q_{f,\Gamma}\right),
\end{gathered}
\label{eq:rh}
\end{equation}
where the subscripts $p$ and $f$ denote the pressurized trace and the liquid
projection of the stratified trace at the front, and $h_{s,\Gamma}$ follows
from the elastic closure~(\ref{eq:elastic}). The pressurized side is
liquid-full, so gas does not cross the front; the moving-frame gas condition
gives
\begin{equation}
u_{g,\Gamma}=w_\Gamma .
\label{eq:gaskin}
\end{equation}
Because the water-hammer speed $a$ is much larger than mixed-flow front
speeds, one pressurized acoustic characteristic always reaches the front,
supplying the compatibility relation
\begin{equation}
u_{p,\Gamma}-u_{p,L}+\frac{g}{a}\left(h_{s,\Gamma}-h_{s,L}\right)=0,
\label{eq:compat}
\end{equation}
with $L$ the adjacent pressurized state. Equations
(\ref{eq:rh})--(\ref{eq:compat}) do not determine all interface unknowns;
the remaining information depends on the position of $w_\Gamma$ relative to
the stratified liquid signal speeds
$s_\pm=u_{l,R}\pm c_{l,R}$, $c_{l,R}=\sqrt{\Lambda_{d,R}A_R}$, which
partition the admissible front speeds into the slow ($w<s_-$), middle
($s_-\le w\le s_+$), and fast ($w>s_+$) active sets.

Because the correct active set is not known before $w_\Gamma$ is computed,
and branch-dependent nonlinear iteration at every front is fragile in
transient computations, the closure of \cite{xue2026model} is a fixed-cost
one-step construction. Substituting the compatibility
relation~(\ref{eq:compat}) and the mass jump into the momentum jump reduces
the closure, for any candidate speed $w$, to a scalar quadratic for
$A_{p,\Gamma}(w)$ and a scalar momentum residual $R(w)$. A finite-area
predictor built from the adjacent states,
\begin{equation}
w_0=\frac{A_{p,L}u_{p,L}-A_R u_{l,R}}{A_{p,L}-A_R},
\label{eq:predictor}
\end{equation}
is projected onto each (slightly shrunken) active-set interval, one
closed-form Newton-type residual correction is applied per interval, and the
accepted candidate is the one minimizing a penalized criterion combining the
Rankine--Hugoniot residual, the distance to its active set, and positivity
penalties. No Bernoulli-type energy relation is imposed at the front. The
selected trace defines the moving-frame flux pair
$\mathbf{G}_{p,\Gamma}=\mathbf{F}_{p,\Gamma}-w_\Gamma\mathbf{U}_{p,\Gamma}$
and
$\mathbf{G}_{f,\Gamma}=\mathbf{F}_{f,\Gamma}-w_\Gamma\mathbf{U}_{f,\Gamma}$
passed to the cut-cell update below.

\subsection{Finite-volume update and cut-cell front tracking}
\label{sec:model:fv}

Away from fronts, both branches are advanced by the same explicit
finite-volume machinery: MUSCL--TVD primitive reconstruction with the minmod
limiter, branch-dependent Rusanov fluxes (scalar dissipation
$\max(|u_p|+a)$ on the pressurized branch; block-diagonal dissipation with
gas speed $|u_g|+c_g$ and stratified liquid speed $|u_l|+c_\Lambda$,
$c_\Lambda=\sqrt{\max(\Lambda_d A_l,0)+c_{\mathrm{num}}^2}$, on the
stratified branch), unsplit cell-centered sources, and SSP--RK2 time
integration. The time step is the minimum of the branch CFL limits and a
front-motion restriction
$\Delta t_\Gamma=\tfrac{1}{2}\min(\Delta x_p,\Delta x_f)/|w_\Gamma^n|$,
which prevents the front from sweeping more than half a cell per step.

The cell hosting a front is divided into pressurized and stratified
sub-cells of lengths $\ell_p$ and $\ell_f$ with
$\ell_p+\ell_f=\Delta x$. At the beginning of each step the adjacent branch
states are passed to the closure of Section~\ref{sec:model:rtrp}; the front
is then advanced as $x_\Gamma^{n+1}=x_\Gamma^{n}+w_\Gamma\Delta t$, and each
sub-cell is updated by its own branch flux at the outer face and the
moving-frame interface flux at the front. Because both sub-cell balances use
the same $w_\Gamma$ and the same interface trace, the interface fluxes and
swept-volume contributions cancel exactly for the conserved liquid
components and gas mass: the front neither creates nor destroys liquid or
gas. When the front crosses a cell face, the affected cells are remapped
conservatively (event-based remap), and a bounded, localized residual
relaxation corrects the remaining interface mismatch after an accepted
step. Trapped gas regions are identified after each step as connected gas
volumes bounded by liquid seals, and their pressures are evaluated from the
isothermal equation of state on the current pocket mass and volume.

\subsection{Vertical-shaft branch}
\label{sec:model:riser}

A vertical shaft (ventilation tower, riser, or dropshaft barrel) is
discretized as a separate one-dimensional finite-volume branch using the same
area--momentum variables as the horizontal conduit, with the pipe inclination
set to $\theta=90^{\circ}$ so that gravity acts axially. The shaft top is
open to the atmosphere; the free surface inside the shaft and any gas core
penetrating below it are represented by the axial profile of the gas volume
fraction $\alpha_g(z)$, not by a prescribed bubble shape. The shaft state
therefore distinguishes the free-surface elevation $Y_{fs}$ (top of the
continuous liquid column) from the air--water interface elevation $Y_{int}$
(nose of the gas core rising through the column), which are exactly the two
trajectories reported in shaft experiments. Geysering is diagnosed, following
the experimental definition of \cite{vasconcelos2011}, when $Y_{fs}$ reaches
the shaft top before $Y_{int}$ catches the free surface.

\subsection{Junction and local-loss closures}
\label{sec:model:junction}

The shaft base is coupled to the horizontal conduit through a junction
closure driven by the local pipe state and the trapped-pocket pressure. The
junction exchange uses the pocket equation-of-state static head as the
driving head, with the following local-loss elements retained from the
frozen configuration used for all runs in Section~\ref{sec:tests}:

\begin{itemize}
\item a junction entry/exit loss coefficient $K_j=0.75$ on the through-flow
between pipe and shaft base;
\item an additional churn loss $K_{g}=8.0$ applied while the junction carries
gas, representing the strongly dissipative counter-current air-up/water-down
exchange (``glugging'') at the shaft mouth;
\item a butterfly-valve loss $K_v=2.0$ at the release valve, whose disc
remains in the flow after opening, with a finite hand-opening time
$t_v=0.25$~s represented by throttling the loss coefficient during opening;
\item a counter-current flooding (CCFL) limit of Wallis type on reverse
liquid drainage at the shaft top while gas is venting, with limiting velocity
scale $0.5\sqrt{gD_t}$, where $D_t$ is the shaft diameter.
\end{itemize}

Three further elements govern the shaft base in the driven regime, i.e.\
when the trapped pocket holds an overpressure margin over the shaft column
($H_{a0}>Y_{fs,0}$) and gas has reached the junction; they are stated here
and motivated with the narrow-tower test in
Section~\ref{sec:tests:vw2011:caseB}: (i) the shaft base is driven by the
pipe \emph{crown} pressure (the elevation at which the shaft physically
taps the pipe), one bore below the invert piezometric line while the
junction is water-sealed; (ii) the shaft gas core connected to the pocket
is pocket-fed---mouth influx raises the core nose, by the same
fill-to-body-level closure used for the crown cavity in the horizontal
pipe---rather than accumulating at the base cell under dispersed-bubble
drag; and (iii) the entry gas flux is bounded by the Wallis counter-current
flooding scale $j_g\le C_w^{2}\sqrt{gD_t\,\Delta\rho/\rho_l}$ with
$C_w=0.9$. These closures use geometric facts and literature-scale
constants; in the undriven regime (Test~A of Campaign~1) they are inactive
and the base retains the invert-datum coupling, a dual treatment whose
consolidation is discussed in Section~\ref{sec:discussion}.

Once a trapped pocket becomes connected to the shaft and gas breaks through
the liquid seal, venting latches on: the pocket continues to discharge
through the shaft and is not re-sealed by intermittent slugs. This latch
reflects the experimental observation that after breakthrough the release
proceeds as a sustained, if oscillatory, discharge rather than a sequence of
independent seal events.

\subsection{Numerical configuration}
\label{sec:model:numerics}

All computations in Section~\ref{sec:tests} use one frozen solver
configuration: horizontal grid spacing $\Delta s=0.02$~m, shaft grid spacing
$\Delta z=0.01$~m, CFL number 0.35, and a reduced numerical elastic wave
speed $a=28$~m/s in the liquid-full branch. The reduced wave speed keeps the
liquid-full reaches effectively incompressible on the pocket and shaft
timescales (the gravity-wave scale is $\sqrt{gD}\approx 1$~m/s) while
avoiding stiff acoustic ringing at grid-scale micro-voids; the geysering
dynamics studied here are pocket-compression and gravity-wave driven, not
water-hammer driven. Gas thermodynamics is isothermal with $R=287$~J/(kg\,K)
and $T=293$~K, and the liquid density is $\rho_l=998$~kg/m$^3$. No loss
coefficient or numerical parameter is adjusted between test cases; the only
case-dependent element is the regime-conditioned shaft-base treatment of
Section~\ref{sec:model:junction} (crown-datum, pocket-fed, flooding-limited
in the driven regime; invert-datum otherwise), which switches on the
physical state, not on the test identity.
